\hypertarget{appendix-7281c15ac44a}{}
\section*{Retinal encoding and receptive fields}
\section{Retinal encoding and receptive fields}\label{retinal-encoding-and-receptive-fields}

\textbf{Curriculum depth:} Introduced · \textbf{First seen:} XR01 · \textbf{Latest development:} XR01

\subsection{Definition and mechanism}\label{definition-and-mechanism}

Rods/cones transduce light into neural signals; they are not passive RGB pixels. Receptive-field models describe how neighboring inputs affect responses. A balanced difference of normalized Gaussians suppresses constants, unlike arbitrary unequal-gain kernels.

\subsection{Prerequisites}\label{prerequisites}

\begin{itemize}
\tightlist
\item
  \hyperref[appendix-215c2570199b]{Light spectra and radiometry} --- The light quantity and units must be established before its interaction or visual effect is modeled.
\end{itemize}

\subsection{Used by}\label{used-by}

\begin{itemize}
\item
  \hyperref[appendix-b491fc48c54b]{Visual acuity and eccentricity}
\item
  \hyperref[appendix-b3aa9e8e01ad]{Contrast sensitivity adaptation and masking}
\item
  \hyperref[appendix-47c51ea026fb]{Learned image features and CNNs} \#\# Recall and next depth
\end{itemize}

Explain why unequal Gaussian gains retain DC. Pending: calibrated physiological models and cortical processing. Written coverage is not a claim of learner mastery or accepted numbered-curriculum progress.

\subsection{Source and navigation}\label{source-and-navigation}

\href{Continuous\%20Notes/University\%20-\%20Worlds,\%20Images\%20and\%20Inference\%20-\%20Graphics,\%20Vision\%20and\%20Photography\#The\%20eye\%20is\%20an\%20optical\%20instrument\%20with\%20neural\%20processing\%20inside\%20it}{Expanded teaching explanation} · \href{}{Independent coverage (01-Foundations-I, appendix-fd090d982a15 § XR01 - Independent seeing and representation source)} · \href{}{Dependency Map (02-Foundations-II, appendix-4ad47bdd269a)}
